\chapter{Inversions and Voice Leading: Moving Fingers Inches, Not Feet}
\label{ch:inversions-and-voice-leading}
\chapterline

Watch a beginner try to play a simple three-chord song --- $C$ major to $F$ major to $G$ major --- using only the root-position shapes from Chapter~3. They press $C - E - G$, stare down in mild alarm, launch their whole arm four keys to the right to grab $F - A - C$, launch again to $G - B - D$, and leap all the way back down to $C - E - G$. Every chord change feels like jumping across stepping stones in the dark, and the steady beat breaks apart.

Now watch an experienced pianist play those exact same three chords. Their right hand barely moves an inch. Their wrist stays quiet, they never even glance down at the keyboard, and the chords melt smoothly into one another like a four-part choir. What are they doing differently? They are folding the chords into \textbf{inversions}.

\section{Three Shapes of Every Chord, One Gap Detector}

Remember from Chapter~1 that shifting any note up or down by $+12$ half steps (one octave, a $2:1$ wave ratio) keeps the \emph{exact same pitch name}. A chord is defined by its set of pitch names, regardless of which note sits on the bottom. Whenever you take the lowest note of a three-note chord and flip it $+12$ half steps up to the top, you create an \textbf{inversion} of that chord.

Follow $C$ major ($C, E, G$) across its three hand shapes in Figure~\ref{fig:three-inversions}:

\begin{enumerate}
  \item \textbf{Root position ($C4 - E4 - G4$, right-hand fingers $\mathbf{1 - 3 - 5}$):} The root $C$ sits at the bottom. Both gaps are small thirds ($1$ skipped white key in each gap), spanning five white keys total.
  \item \textbf{First inversion ($E4 - G4 - C5$, right-hand fingers $\mathbf{1 - 2 - 5}$):} Flip bottom $C4$ up one octave to $C5$. Now $E4 - G4$ is a snug third at the bottom ($1$ skipped white key), while $G4 - C5$ opens into a wider \textbf{fourth} on top ($2$ skipped white keys, $5$ half steps). Notice the right-hand fingering $\mathbf{1 - 2 - 5}$: placing your index finger ($\mathbf{2}$) on $G4$ leaves a relaxed, natural stretch between finger~$2$ and finger~$5$ across the wider top gap.
  \item \textbf{Second inversion ($G4 - C5 - E5$, right-hand fingers $\mathbf{1 - 3 - 5}$):} Flip bottom $E4$ up one octave to $E5$. Now the wider fourth ($G4 - C5$, $2$ skipped white keys) sits at the bottom --- right where the flexible web between your thumb ($\mathbf{1}$) and middle finger ($\mathbf{3}$) opens effortlessly --- and the snug third ($C5 - E5$) sits on top. Flip $G4$ up one more octave and you return to root position ($C5 - E5 - G5$).
\end{enumerate}

\begin{formalbox}
Every three-note triad $\set{r, t, f}$ has exactly three close-position voicings within an octave:
\begin{itemize}
  \item \textbf{Root position:} $(r,\; t,\; f)$ --- spans $7$ half steps (a fifth); right-hand fingering $\mathbf{1 - 3 - 5}$ (left hand $\mathbf{5 - 3 - 1}$).
  \item \textbf{First inversion:} $(t,\; f,\; r + 12)$ --- spans $8$ or $9$ half steps (a sixth), with a $5$-semitone fourth on \textbf{top}; right-hand fingering $\mathbf{1 - 2 - 5}$ (left hand $\mathbf{5 - 3 - 1}$).
  \item \textbf{Second inversion:} $(f,\; r + 12,\; t + 12)$ --- spans $8$ or $9$ half steps (a sixth), with a $5$-semitone fourth on \textbf{bottom}; right-hand fingering $\mathbf{1 - 3 - 5}$ (left hand $\mathbf{5 - 2 - 1}$).
\end{itemize}
\end{formalbox}

\subsection*{The Gap Detector: name any inverted chord in one glance}

Suppose your fingers land on an unfamiliar inverted triad such as $A - D - F\sharp$, and you want to know its root name right away without flipping notes in your head. Use the \textbf{Gap Detector}:
\begin{quote}
\emph{Look for the wider gap ($5$ half steps, skipping two scale steps instead of one). The note sitting immediately \textbf{above} that wider gap is \textbf{always the root} of the chord.}
\end{quote}
Check it on $A - D - F\sharp$: between $A$ and $D$ sits a wide $5$-semitone fourth ($B$ and $C\sharp$ skipped), whereas $D$ to $F\sharp$ is a snug $4$-semitone third. Which note sits directly \emph{above} the wide gap? \textbf{$D$} --- so $A - D - F\sharp$ is a \textbf{$D$ major chord} (in second inversion). Check $E - G - C$: the wide gap sits between $G$ and $C$, and the note right above that gap is \textbf{$C$} ($C$ major, first inversion). If neither gap is wider, the chord is already in root position and the bottom note is the root.

\begin{figure}[htbp]
\centering
\begin{tikzpicture}
  % Three cleanly stacked strips on C4..G5 (12 white keys, wk = 0.80cm -> 9.60cm wide)
  \def\wk{0.80}
  \def\hk{1.75}

  % Row 1: Root Position C4 - E4 - G4
  \begin{scope}[shift={(2.65,6.20)}]
    \node[anchor=east, font=\footnotesize\bfseries, text=sourceblue, align=right] at (-0.25,0.88)
      {Root Position\\$C4 - E4 - G4$\\ {\scriptsize RH: $1 - 3 - 5$}};
    \foreach \k in {0,...,11} {
      \filldraw[fill=white, draw=black!65, line width=0.55pt] ({\k*\wk},0) rectangle ({(\k+1)*\wk},\hk);
    }
    \foreach \k in {0,2,4} {
      \filldraw[fill=sourceblue!18, draw=black!65, line width=0.55pt] ({\k*\wk},0) rectangle ({(\k+1)*\wk},\hk);
    }
    \foreach \k in {0,1,3,4,5,7,8,10} {
      \filldraw[fill=black!85, draw=black, line width=0.5pt] ({(\k+1)*\wk-0.21},0.72) rectangle ({(\k+1)*\wk+0.21},\hk);
    }
    \foreach \k/\fn/\nm in {0/1/C4, 2/3/E4, 4/5/G4} {
      \fill[sourceblue] ({(\k+0.5)*\wk},0.42) circle (0.10);
      \node[font=\scriptsize\bfseries, text=sourceblue] at ({(\k+0.5)*\wk},-0.26) {\nm\ (\fn)};
    }
    \draw[<->, thick, >=Latex, quotientorange!85!black, dashed, shorten <=2pt, shorten >=2pt]
      ({0.5*\wk},1.83) to[out=22,in=158] node[above=2pt, font=\scriptsize\bfseries, text=quotientorange!85!black] {Flip bottom $C4$ up $+12$ semitones to $C5$} ({7.5*\wk},1.83);
  \end{scope}

  % Row 2: First Inversion E4 - G4 - C5
  \begin{scope}[shift={(2.65,3.10)}]
    \node[anchor=east, font=\footnotesize\bfseries, text=targetgreen!80!black, align=right] at (-0.25,0.88)
      {1st Inversion\\$E4 - G4 - C5$\\ {\scriptsize RH: $1 - 2 - 5$}};
    \foreach \k in {0,...,11} {
      \filldraw[fill=white, draw=black!65, line width=0.55pt] ({\k*\wk},0) rectangle ({(\k+1)*\wk},\hk);
    }
    \foreach \k in {2,4} {
      \filldraw[fill=targetgreen!16, draw=black!65, line width=0.55pt] ({\k*\wk},0) rectangle ({(\k+1)*\wk},\hk);
    }
    \filldraw[fill=sourceblue!22, draw=black!65, line width=0.55pt] ({7*\wk},0) rectangle ({8*\wk},\hk);
    \foreach \k in {0,1,3,4,5,7,8,10} {
      \filldraw[fill=black!85, draw=black, line width=0.5pt] ({(\k+1)*\wk-0.21},0.72) rectangle ({(\k+1)*\wk+0.21},\hk);
    }
    \fill[targetgreen!80!black] ({2.5*\wk},0.42) circle (0.10);
    \fill[targetgreen!80!black] ({4.5*\wk},0.42) circle (0.10);
    \fill[sourceblue] ({7.5*\wk},0.42) circle (0.11);
    \node[font=\scriptsize\bfseries, text=targetgreen!80!black] at ({2.5*\wk},-0.26) {E4 (1)};
    \node[font=\scriptsize\bfseries, text=targetgreen!80!black] at ({4.5*\wk},-0.26) {G4 (2)};
    \node[font=\scriptsize\bfseries, text=sourceblue] at ({7.5*\wk},-0.26) {C5 (5)};
    % Wide gap bracket on top
    \draw[decorate, decoration={brace, amplitude=4pt}, thick, sourceblue]
      ({4.55*\wk},1.83) -- ({7.45*\wk},1.83)
      node[midway, above=3pt, font=\scriptsize\bfseries, text=sourceblue] {Wider gap (4th) $\to$ Root $C5$ on top};
  \end{scope}

  % Row 3: Second Inversion G4 - C5 - E5
  \begin{scope}[shift={(2.65,0.00)}]
    \node[anchor=east, font=\footnotesize\bfseries, text=structurepurple!85!black, align=right] at (-0.25,0.88)
      {2nd Inversion\\$G4 - C5 - E5$\\ {\scriptsize RH: $1 - 3 - 5$}};
    \foreach \k in {0,...,11} {
      \filldraw[fill=white, draw=black!65, line width=0.55pt] ({\k*\wk},0) rectangle ({(\k+1)*\wk},\hk);
    }
    \filldraw[fill=structurepurple!16, draw=black!65, line width=0.55pt] ({4*\wk},0) rectangle ({5*\wk},\hk);
    \filldraw[fill=sourceblue!22, draw=black!65, line width=0.55pt] ({7*\wk},0) rectangle ({8*\wk},\hk);
    \filldraw[fill=structurepurple!16, draw=black!65, line width=0.55pt] ({9*\wk},0) rectangle ({10*\wk},\hk);
    \foreach \k in {0,1,3,4,5,7,8,10} {
      \filldraw[fill=black!85, draw=black, line width=0.5pt] ({(\k+1)*\wk-0.21},0.72) rectangle ({(\k+1)*\wk+0.21},\hk);
    }
    \fill[structurepurple!85!black] ({4.5*\wk},0.42) circle (0.10);
    \fill[sourceblue] ({7.5*\wk},0.42) circle (0.11);
    \fill[structurepurple!85!black] ({9.5*\wk},0.42) circle (0.10);
    \node[font=\scriptsize\bfseries, text=structurepurple!85!black] at ({4.5*\wk},-0.26) {G4 (1)};
    \node[font=\scriptsize\bfseries, text=sourceblue] at ({7.5*\wk},-0.26) {C5 (3)};
    \node[font=\scriptsize\bfseries, text=structurepurple!85!black] at ({9.5*\wk},-0.26) {E5 (5)};
    % Wide gap bracket on bottom-left
    \draw[decorate, decoration={brace, amplitude=4pt}, thick, sourceblue]
      ({4.55*\wk},1.83) -- ({7.45*\wk},1.83)
      node[midway, above=3pt, font=\scriptsize\bfseries, text=sourceblue] {Wider gap (4th) $\to$ Root $C5$ right above gap};
  \end{scope}
\end{tikzpicture}
\caption{Folding $C$ major across two octaves ($C4$ to $G5$). Flipping the bottom note up $+12$ semitones moves from root position ($1-3-5$) to first inversion ($1-2-5$) to second inversion ($1-3-5$). In every inversion, the root ($C5$, shaded blue) sits immediately above the wider fourth gap.}
\label{fig:three-inversions}
\end{figure}

\section{Smooth Voice Leading: The Anchor-Note Law}

Think of the three notes in your right-hand chord as three singers standing side by side. No real singer leaps five scale steps up and down on every bar when a note one step away belongs to the next chord. Moving each individual note along the shortest possible path from chord to chord is called \textbf{smooth voice leading}, and one physical habit makes it automatic at the piano:

\begin{intuitionbox}
\textbf{The Anchor-Note Law:} When moving from one chord to the next, \emph{keep any shared note frozen under the finger already holding it}, and step your remaining fingers by just one key to the nearest notes of the new chord.
\end{intuitionbox}

Walk through a complete song progression around Middle~$C$ and watch Figure~\ref{fig:voice-leading-comparison} compare clunky root-position jumps against smooth anchored inversions:

\begin{enumerate}
  \item \textbf{Start on $C$ major ($C4 - E4 - G4$, root position, fingers $1 - 3 - 5$).}
  \item \textbf{Move $C \to F$ major ($F, A, C$):} Which note do $C$ major and $F$ major share? \textbf{$C4$}. Freeze your thumb ($\mathbf{1}$) on $C4$, and step fingers $\mathbf{3}$ and $\mathbf{5}$ up just \emph{one white key} from $E4 - G4$ to $F4 - A4$. You are now holding $C4 - F4 - A4$ ($F$ major, second inversion, fingers $1 - 3 - 5$). Total arm jump: zero inches.
  \item \textbf{Move $C \to G$ major ($G, B, D$):} Return to $C4 - E4 - G4$. Which note do $C$ major and $G$ major share? \textbf{$G4$}. Freeze your pinky ($\mathbf{5}$) on $G4$, and step your bottom two notes down just \emph{one white key} onto $B3 - D4$ using fingers $\mathbf{1 - 2 - 5}$. You are now holding $B3 - D4 - G4$ ($G$ major, first inversion).
  \item \textbf{Move $G \to Am \to F \to C$:} From $G$ ($B3 - D4 - G4$), step all three voices smoothly up one white key to $C4 - E4 - A4$ ($A$ minor, first inversion, $1 - 2 - 5$). Now move from $Am$ ($C4 - E4 - A4$) to $F$ ($C4 - F4 - A4$): \emph{two} shared notes ($C4$ and $A4$) stay completely still while the middle note moves a single half step ($E4 \to F4$). Finally step $F4 - A4$ down one white key to $E4 - G4$ and you are back home on $C$.
\end{enumerate}

\begin{figure}[htbp]
\centering
\begin{tikzpicture}
  % Left panel (x = 0..5.8): Clunky Root-Position Jumps
  \begin{scope}[shift={(0,0)}]
    \node[anchor=south, font=\footnotesize\bfseries, text=trapred] at (2.90,4.55)
      {Panel (a): Clunky Root-Position Leaps};

    \node[draw=trapred!75, fill=trapred!7, rounded corners=3pt, inner xsep=7pt, inner ysep=5pt, minimum width=5.40cm, font=\footnotesize] (r1) at (2.90,3.85)
      {\textbf{$C$ (Root):} $C4 \;-\; E4 \;-\; G4$};
    \node[font=\scriptsize\bfseries, text=trapred, inner sep=2pt] (rl1) at (2.90,2.95)
      {Hand leaps $+5$ semitones};
    \node[draw=trapred!75, fill=trapred!7, rounded corners=3pt, inner xsep=7pt, inner ysep=5pt, minimum width=5.40cm, font=\footnotesize] (r2) at (2.90,2.05)
      {\textbf{$F$ (Root):} $\qquad\; F4 \;-\; A4 \;-\; C5$};
    \node[font=\scriptsize\bfseries, text=trapred, inner sep=2pt] (rl2) at (2.90,1.15)
      {Leaps $+2$ more, then $-7$ back};
    \node[draw=trapred!75, fill=trapred!7, rounded corners=3pt, inner xsep=7pt, inner ysep=5pt, minimum width=5.40cm, font=\footnotesize] (r3) at (2.90,0.25)
      {\textbf{$G$ (Root):} $\qquad\quad\; G4 \;-\; B4 \;-\; D5$};

    \draw[thick, trapred, shorten <=1pt, shorten >=1pt] (r1.south) -- (rl1.north);
    \draw[->, thick, >=Latex, trapred, shorten <=1pt, shorten >=1pt] (rl1.south) -- (r2.north);
    \draw[thick, trapred, shorten <=1pt, shorten >=1pt] (r2.south) -- (rl2.north);
    \draw[->, thick, >=Latex, trapred, shorten <=1pt, shorten >=1pt] (rl2.south) -- (r3.north);
  \end{scope}

  % Divider
  \draw[gray!40, dashed, thick] (6.35,-0.15) -- (6.35,4.85);

  % Right panel (x = 6.90..13.0): Smooth Voice-Led Inversions in one 6-key zone B3..A4!
  \begin{scope}[shift={(6.90,0)}]
    \node[anchor=south, font=\footnotesize\bfseries, text=targetgreen!80!black] at (3.05,4.55)
      {Panel (b): Smooth Anchored Inversions};

    \node[draw=sourceblue!80, fill=sourceblue!7, rounded corners=3pt, inner xsep=7pt, inner ysep=5pt, minimum width=5.80cm, font=\footnotesize] (s1) at (3.05,3.85)
      {\textbf{$C$ (Root):} $\;\; \textcolor{targetgreen!80!black}{\mathbf{C4}} - E4 - \textcolor{targetgreen!80!black}{\mathbf{G4}}$ \hfill {\scriptsize $(1\text{-}3\text{-}5)$}};
    \node[font=\scriptsize\bfseries, text=targetgreen!80!black, inner sep=2pt] (sl1) at (3.05,2.95)
      {$C4$ anchored; top two step $+1$ key};
    \node[draw=targetgreen!80!black, fill=targetgreen!8, rounded corners=3pt, inner xsep=7pt, inner ysep=5pt, minimum width=5.80cm, font=\footnotesize] (s2) at (3.05,2.05)
      {\textbf{$F$ (2nd Inv):} $\textcolor{targetgreen!80!black}{\mathbf{C4}} - F4 - A4$ \hfill {\scriptsize $(1\text{-}3\text{-}5)$}};
    \node[font=\scriptsize\bfseries, text=sourceblue, inner sep=2pt] (sl2) at (3.05,1.15)
      {Step smoothly into $B3 - D4 - G4$};
    \node[draw=quotientorange!85!black, fill=quotientorange!8, rounded corners=3pt, inner xsep=7pt, inner ysep=5pt, minimum width=5.80cm, font=\footnotesize] (s3) at (3.05,0.25)
      {\textbf{$G$ (1st Inv):} $B3 - D4 - \textcolor{targetgreen!80!black}{\mathbf{G4}}$ \hfill {\scriptsize $(1\text{-}2\text{-}5)$}};

    \draw[thick, targetgreen!80!black, shorten <=1pt, shorten >=1pt] (s1.south) -- (sl1.north);
    \draw[->, thick, >=Latex, targetgreen!80!black, shorten <=1pt, shorten >=1pt] (sl1.south) -- (s2.north);
    \draw[thick, sourceblue, shorten <=1pt, shorten >=1pt] (s2.south) -- (sl2.north);
    \draw[->, thick, >=Latex, sourceblue, shorten <=1pt, shorten >=1pt] (sl2.south) -- (s3.north);
  \end{scope}
\end{tikzpicture}
\caption{Left: Jumping between root-position triads forces your right hand to leap across a nine-key span ($C4$ to $D5$), breaking eye contact with the music. Right: Using anchored inversions keeps all three chords inside one calm six-key zone ($B3$ to $A4$) around Middle~$C$.}
\label{fig:voice-leading-comparison}
\end{figure}

\section{Two-Hand Piano Accompaniment That Sounds Full Today}

If your right hand is playing inverted chords such as $C4 - F4 - A4$ (where $C$, not $F$, sits at the bottom of your right hand), why does the listener still hear a clear $F$ major chord rather than a blurred cluster? Because your \textbf{left hand} plays a low $F2$ or $F3$ in the bass, establishing the true root for the whole room.

\begin{trapbox}
Beginners often try to play the exact same closed three-note triad in \emph{both} hands at once --- putting $C2 - E2 - G2$ down in the low bass alongside $C4 - E4 - G4$ above. Try pressing $C2 - E2 - G2$ low on a real piano and listen closely: it sounds thick and muddy. Thick low piano strings produce strong overtones that clash when thirds sit closer than an octave below $C3$ (pianists call this the \textbf{low-interval mud zone}). Keep closed three-note chords above $C3$, and give the low bass wide, open intervals: a single bass root, a fifth ($C2 - G2$), or an octave ($C2 - C3$).
\end{trapbox}

Figure~\ref{fig:two-hand-layout} shows the clean studio spacing used by ballad, pop, gospel, and classical pianists alike: left hand anchors wide roots or fifths down in the $C2$--$C3$ bass register, while right hand voices smooth three-note triads right around Middle~$C4$.

\begin{figure}[htbp]
\centering
\begin{tikzpicture}
  % 22 white keys from C2 to C5: wk = 0.62cm -> 22 * 0.62 = 13.64cm wide!
  \def\wk{0.62}
  \def\hk{2.35}

  \foreach \k in {0,...,21} {
    \filldraw[fill=white, draw=black!65, line width=0.5pt] ({\k*\wk},0) rectangle ({(\k+1)*\wk},\hk);
  }
  % Light mud-zone warning tint on D2..B2 (k=1..6) except C2(0) and G2(4) which hold our open 5th!
  \filldraw[fill=sourceblue!20, draw=black!65, line width=0.5pt] ({0*\wk},0) rectangle ({1*\wk},\hk);
  \filldraw[fill=sourceblue!20, draw=black!65, line width=0.5pt] ({4*\wk},0) rectangle ({5*\wk},\hk);

  % Right-hand voice-led triad around Middle C4 (C4=14, E4=16, G4=18)
  \foreach \k in {14,16,18} {
    \filldraw[fill=targetgreen!20, draw=black!65, line width=0.5pt] ({\k*\wk},0) rectangle ({(\k+1)*\wk},\hk);
  }

  % Black keys across 3 octaves:
  % Oct 2 (0..6): 0,1, 3,4,5
  % Oct 3 (7..13): 7,8, 10,11,12
  % Oct 4 (14..20): 14,15, 17,18,19
  \foreach \k in {0,1,3,4,5, 7,8,10,11,12, 14,15,17,18,19} {
    \filldraw[fill=black!85, draw=black, line width=0.45pt] ({(\k+1)*\wk-0.18},0.94) rectangle ({(\k+1)*\wk+0.18},\hk);
  }

  % Dots on LH C2(0), G2(4) and RH C4(14), E4(16), G4(18)
  \fill[sourceblue] ({0.5*\wk},0.44) circle (0.11);
  \fill[sourceblue] ({4.5*\wk},0.44) circle (0.11);
  \foreach \k in {14,16,18} { \fill[targetgreen!80!black] ({(\k+0.5)*\wk},0.44) circle (0.11); }

  % Octave landmark labels below keys
  \node[font=\scriptsize\bfseries, text=sourceblue] at ({0.5*\wk},-0.26) {C2};
  \node[font=\scriptsize\bfseries, text=sourceblue] at ({4.5*\wk},-0.26) {G2};
  \node[font=\scriptsize, text=gray!80!black] at ({7.5*\wk},-0.26) {C3};
  \node[font=\scriptsize\bfseries, text=targetgreen!80!black] at ({14.5*\wk},-0.26) {C4 (Mid C)};
  \node[font=\scriptsize, text=gray!80!black] at ({21.5*\wk},-0.26) {C5};

  % Clean brackets above LH and RH zones + middle spacing callout
  \draw[decorate, decoration={brace, amplitude=5pt}, thick, sourceblue]
    ({0.1*\wk},2.46) -- ({4.9*\wk},2.46)
    node[midway, above=5pt, font=\scriptsize\bfseries, text=sourceblue, align=center] {Left Hand: Open Bass\\Root ($C2$) or Fifth ($C2$--$G2$)};

  \node[draw=trapred!70, fill=trapred!6, rounded corners=3pt, inner xsep=6pt, inner ysep=4pt, font=\scriptsize, align=center] at ({9.5*\wk},2.92)
    {Avoid closed 3-note triads\\below $C3$ (muddy overtones)};

  \draw[decorate, decoration={brace, amplitude=5pt}, thick, targetgreen!80!black]
    ({14.0*\wk},2.46) -- ({19.0*\wk},2.46)
    node[midway, above=5pt, font=\scriptsize\bfseries, text=targetgreen!80!black, align=center] {Right Hand: Voice-Led\\Triads Near Middle $C4$};
\end{tikzpicture}
\caption{Balanced two-hand spacing across three octaves ($C2$ to $C5$). Left hand plays clean single roots, fifths, or octaves in the bass while right hand keeps three-note inversions inside a warm zone near Middle~$C4$.}
\label{fig:two-hand-layout}
\end{figure}

Once your hands sit in those two zones, three rhythmic accompaniment patterns bring any chord sheet to life immediately:
\begin{enumerate}
  \item \textbf{The Steady Pulse (Singer-Songwriter Quarter Groove):} Hold a single bass note in your left hand for four steady counts while your right hand presses its three-note chord gently on every beat ($1, 2, 3, 4$).
  \item \textbf{The Rolling Wave (Broken-Chord Ballad):} Hold the left-hand bass note on beat~$1$, while your right hand unpacks its three notes one finger at a time: $\text{Bottom} \to \text{Middle} \to \text{Top} \to \text{Middle}$ across the bar.
  \item \textbf{The Pop Bounce (Together -- Right -- Left -- Right):} Press left-hand pinky ($C2$) and right-hand chord \emph{together} on beat~$1$, right-hand chord alone on beat~$2$, left-hand thumb ($G2$) on beat~$3$, and right-hand chord on beat~$4$.
\end{enumerate}

\section{Slash Chords Demystified: Walking Basslines Made Easy}

Open almost any songbook or guitar-and-piano chord sheet and you will meet symbols with a forward slash, such as \textbf{$C/E$}, \textbf{$G/B$}, or \textbf{$D/F\sharp$}. Beginners sometimes mistake these for fractions or two chords mashed together. The rule could not be simpler:
\[
  \underbrace{\text{Chord Name}}_{\text{Right-hand triad}} \;\Big/\; \underbrace{\text{Bass Note}}_{\text{Left-hand single note}}.
\]
Thus \textbf{$G/B$} literally reads ``\emph{$G$ major chord over a $B$ bass note}'': your right hand plays any smooth $G$ major chord around Middle~$C$, while your left hand presses the single note $B$ down in the bass.

Why do songwriters write $G/B$ instead of plain $G$? Look at the bass notes across the classic opening line $C \to G/B \to Am \to Am/G \to F$:
\[
  \text{Left-hand bass:} \qquad C3 \;\to\; B2 \;\to\; A2 \;\to\; G2 \;\to\; F2.
\]
Instead of hopping wildly from $C$ down to $G$ and up to $A$, the left hand walks smoothly step by step down the white keys like a cello player, turning five ordinary chords into an unforgettable melodic bassline.

\begin{labbox}
\textbf{Eyes-closed progression drill:} Place your left hand over $C3, B2, A2, G2, F2$ (fingers $1 - 2 - 3 - 4 - 5$ starting with thumb on $C3$) and your right hand on $C4 - E4 - G4$. Play four slow counts per chord across:
\[
  C \;\to\; G/B \;\to\; Am \;\to\; F \;\to\; C,
\]
using smooth right-hand inversions ($C4\text{--}E4\text{--}G4 \to B3\text{--}D4\text{--}G4 \to C4\text{--}E4\text{--}A4 \to C4\text{--}F4\text{--}A4 \to C4\text{--}E4\text{--}G4$). After two slow rounds looking at your fingers, close your eyes and feel every shared anchor note guide your hand without looking.
\end{labbox}

\section*{Exercises}

\begin{exercise}
Spell out all three positions (root position, first inversion, and second inversion) of (a) $F$ major ($F, A, C$) and (b) $G$ major ($G, B, D$). List the recommended right-hand fingering ($1-3-5$ or $1-2-5$) for each shape.
\end{exercise}

\begin{exercise}
Use the Gap Detector (find the wider $5$-semitone fourth gap; the note immediately above it is the root) to identify the chord name and inversion of: (a) $F\sharp4 - A4 - D5$, (b) $E\flat4 - A\flat4 - C5$, and (c) $B\flat3 - E\flat4 - G4$.
\end{exercise}

\begin{exercise}
Starting from root-position $G$ major ($G4 - B4 - D5$), apply the Anchor-Note Law to find the closest smooth right-hand inversions for the four-chord progression $G \to D \to Em \to C$ ($I - V - vi - IV$ in $G$ major). Write down the exact three notes and shared anchor note at each step.
\end{exercise}

\begin{exercise}
What should your left hand and right hand each play when a song sheet shows (a) $D/F\sharp$ between $G$ and $Em$, and (b) $C/G$ right before $G7 \to C$? Trace the three left-hand bass notes of $G \to D/F\sharp \to Em$ to see why songwriters love $D/F\sharp$.
\end{exercise}
